% GENERATED by scripts/manual_tables.py from the adapter registry. Do not edit:
% rerun the script. Hand-written text: blocks/preprocessing/beads/
\ubblock{preprocessing.beads}%
  {BEADS baseline}% what the app calls it
  {BEADS baseline removal (Ning, Selesnick \& Duval 2014)}% registry display name
  {preprocess}% insert-modal tab
  {Signal}% input
  {Signal}% output
  {Adapters/preprocessing\_beads.py}% code
  {preprocessing/beads}% hand-written folder under blocks/
\ubblockdescription{preprocessing.beads}{Removes slow baseline drift without a rolling window: BEADS splits the trace into a smooth baseline, a sparse layer of events and noise, and passes on the trace minus the baseline --- so big drops are not dragged into the baseline and leave no overshoot beside them. Tune `baseline\_period\_s` first: drift slower than it is removed, events faster than it are kept. Set `polarity` to the way your events point. On a span made mostly of events the baseline settles partway into them: event shapes and depths are kept, absolute levels are not reliable there.}
\ubblockflags{preprocessing.beads}{%
  \ubflag{estimate}{}%
}
\ubblocksideinputs{preprocessing.beads}{%
}
\ubblockparams{preprocessing.beads}{%
  \ubparam{baseline\_\allowbreak{}period\_\allowbreak{}s}{7200.0}{number from 60 to 2592000}{Anything slower than this period (seconds) is baseline drift and is removed. Set it above the slowest event you want to keep --- an event as slow as the period is removed with the drift.}%
  \ubparam{polarity}{\texttt{drop}}{one of \texttt{drop}, \texttt{spike}, \texttt{both}}{Which way the events point: drop (troughs), spike (peaks), or both (no preference).}%
  \ubparam{sparsity}{1.0}{number from 0.05 to 50}{How rare events must be (unit-free). Higher: fewer, larger events, more left to the baseline and the noise; lower: more of the trace counted as events.}%
  \ubparam{filter\_\allowbreak{}order}{2}{one of 2}{Sharpness of the drift / event split. Only 2 is offered: order 1 let 7 \% of a drift three times slower than the period through as signal, and 15 \% of one twice as slow.}%
}
